% ---------------------------------------------------------------------------
% Parte 3 — requer no preâmbulo: \usepackage{booktabs,graphicx}
% Figura: copiar code/results/part3_map_vs_dim.png para figures/
% ---------------------------------------------------------------------------
\section{Vocabulários grandes com PCA vs.\ vocabulários pequenos}
\label{sec:part3}

\paragraph{Protocolo.}
Avaliámos BoVW com $K \in \{32,\dots,1024\}$ e VLAD/FV com $K \in \{16,32,64\}$,
sem PCA e com PCA para $d \in \{32,\dots,512\}$ ($d < $ dimensão original e
$d \le N_{\text{treino}}-1$). Para que só $K$ e $d$ variem, fixámos uma
normalização por encoding: BoVW global com Hellinger; VLAD e FV com tiles
$2\times2$ e SSR + L2. Não reutilizámos a escolha da Parte~2, que dava
normalizações diferentes à VLAD consoante $K$. O PCA é ajustado só nas
features de treino (reajustado em cada uma das 3 seeds), e a sua saída é
L2-normalizada. Reportamos o mAP de test de todas as condições (média
$\pm$ desvio).

\begin{table*}[t]
\centering
\caption{Condições A--D (BoVW) e o equivalente para VLAD/FV com $K=64$.
Extração = SIFT + encoding por imagem. Inferência = projeção PCA + SVM por
imagem. O ajuste do PCA demora 0.09\,s (BoVW), 0.84\,s (VLAD) e 1.69\,s (FV).
Accuracy (threshold 0) entre 0.83 e 0.89 em todas as linhas.}
\label{tab:part3}
\small
\setlength{\tabcolsep}{4pt}
\resizebox{\linewidth}{!}{%
\begin{tabular}{llrrcrrrrrr}
\toprule
 & Representação & $d$ & Removido & mAP & F1$_\text{macro}$ & Codebook (s) & Extração (ms) & Treino SVM (s) & Inferência ($\mu$s) & KB/img \\
\midrule
A & BoVW $K{=}256$            &   256 & --      & $0.631 \pm .004$ & 0.699 & 31.1 & 180 & 0.018 & 0.06 & 1.0 \\
B & BoVW $K{=}1024$           &  1024 & --      & $0.682 \pm .003$ & 0.755 & 87.1 & 304 & 0.073 & 0.18 & 4.0 \\
C & BoVW $K{=}1024$ + PCA     &   128 & 87.5\,\% & $0.654 \pm .002$ & 0.747 & 87.1 & 304 & 0.009 & 2.47 & 0.5 \\
C & BoVW $K{=}1024$ + PCA     &    32 & 96.9\,\% & $0.633 \pm .004$ & 0.714 & 87.1 & 304 & 0.003 & 1.23 & 0.125 \\
D & BoVW $K{=}128$            &   128 & --      & $0.590 \pm .004$ & 0.644 & 19.1 & 159 & 0.010 & 0.03 & 0.5 \\
D & BoVW $K{=}32$             &    32 & --      & $0.502 \pm .001$ & 0.601 & 28.2 & 157 & 0.004 & 0.01 & 0.125 \\
\midrule
B & VLAD $K{=}64$             & 20480 & --      & $0.765 \pm .004$ & 0.773 &  8.1 & 166 & 1.369 & 3.32 & 80 \\
C & VLAD $K{=}64$ + PCA       &   128 & 99.4\,\% & $0.730 \pm .002$ & 0.780 &  8.1 & 166 & 0.010 & 40.4 & 0.5 \\
B & FV $K{=}64$               & 41216 & --      & $0.797 \pm .003$ & 0.785 & 55.7 & 234 & 3.074 & 9.59 & 161 \\
C & FV $K{=}64$ + PCA         &   128 & 99.7\,\% & $0.756 \pm .008$ & 0.791 & 55.7 & 234 & 0.010 & 84.6 & 0.5 \\
\bottomrule
\end{tabular}}
\end{table*}

\begin{figure*}[t]
\centering
\includegraphics[width=\linewidth]{figures/part3_map_vs_dim.png}
\caption{mAP de test em função da dimensão final (3 seeds, média $\pm$ desvio).
Cada curva é um codebook reduzido por PCA; o ponto mais à direita é o codebook
sem PCA. No BoVW, a linha tracejada é a condição D ($K=d$, sem PCA).}
\label{fig:part3}
\end{figure*}

\paragraph{Um codebook maior ajuda antes do PCA?}
Sim, mas os ganhos diminuem. O mAP do BoVW passa de 0.502 ($K=32$) para 0.631
($K=256$) e 0.682 ($K=1024$), e a última duplicação já só acrescenta $+0.016$.
O erro de quantização relativo desce de 0.52 para 0.43 ($K=128\to1024$) e é
quase igual em treino e test, por isso o vocabulário maior descreve melhor os
patches sem sobreajustar. O ganho de $K=256$ para 1024 é desigual entre
classes: motorbike $+0.116$ e horse $+0.079$, mas aeroplane só $+0.007$. As
duas primeiras são definidas por estruturas locais finas (rodas e quadro;
pernas e crina), que um vocabulário grosseiro funde com patches de outros
objetos. O aeroplane, com tantos positivos como o motorbike, é plausivelmente
reconhecido pelo contexto (céu, pista), que poucas palavras já captam. VLAD
(0.705$\to$0.765) e FV (0.766$\to$0.797) seguem a mesma tendência
\cite{chatfield2011}.

\paragraph{O PCA retém o ganho?}
Em grande parte, e a perda é gradual (Fig.~\ref{fig:part3}). O BoVW $K=1024$
reduzido a $d=512/256/128$ dá 0.669/0.660/0.654 e retém 75/57/45\,\% da
vantagem sobre $K=256$. Com $d=32$ (97\,\% removido) ainda empata com
$K=256$ completo. Aceitando $-0.02$ de mAP, o PCA remove $\sim$50\,\% das
dimensões no BoVW, 95\,\% na VLAD e 98.8\,\% no FV ($K=64$: 41\,216$\to$512,
0.797$\to$0.783). O PCA nunca melhorou o mesmo codebook para lá do ruído entre
seeds: comprime, mas não regulariza, porque guarda as direções de maior
variância e não as mais discriminativas. Os ganhos reportados exigem whitening
\cite{jegou2012negative}, que não aplicámos. Há uma exceção por explicar: na
VLAD $K=32/64$ o mAP desce em $d=256$ e recupera em $d=512$ ($K=64$:
0.730$\to$0.715$\to$0.742). Não se deve à convergência do SVM nem à
normalização L2 depois do PCA.

\paragraph{À mesma dimensão, grande + PCA supera pequeno?}
Sim, para todo o $d$, e a vantagem cresce com a compressão. C$-$D vale
$+0.131$ em $d=32$, $+0.064$ em $d=128$ e $+0.003$ em $d=512$, este último
dentro do desvio. Na Fig.~\ref{fig:part3}, todas as curvas ficam acima da
tracejada. Um k-means pequeno funde descritores só pela distância no espaço
SIFT. O PCA sobre um histograma grande funde as palavras que co-ocorrem nas
imagens, uma compressão aprendida dos dados que perde menos por dimensão
\cite{jegou2012aggregating}.

\paragraph{Compromisso entre mAP, memória e tempo.}
O PCA atua sobre os custos que já eram pequenos (Tab.~\ref{tab:part3}).
\emph{Construção}: a extração domina (125\,ms de SIFT + até 180\,ms de encoding
por imagem; codebook de 31\,s para 87\,s), e o PCA não a reduz, porque a
representação completa tem de ser calculada primeiro. \emph{Memória e treino}:
é aqui que o PCA ganha. O FV $K=64$ passa de 161 para 0.5\,KB por imagem
($322\times$) e o treino do SVM de 3.1 para 0.01\,s, por $-0.041$ de mAP.
\emph{Inferência}: fica mais lenta (FV: 9.6$\to$85\,$\mu$s), porque projetar
custa $D\cdot d$ operações; só compensa com a projeção guardada, como num
índice. À mesma memória, VLAD/FV com PCA dominam o BoVW: VLAD $K=64$ + PCA~128
(0.5\,KB, 0.730) supera o BoVW $K=1024$ completo (4\,KB, 0.682). A accuracy
quase não varia (0.83--0.89, dominada pelo background); mAP e macro-F1
distinguem as condições.

\paragraph{BoVW vs.\ VLAD/FV.}
A conclusão mantém-se (em $d=128$, de $K=16$ para 64: VLAD 0.682$\to$0.730,
FV 0.720$\to$0.756), e a ordem FV $>$ VLAD $>$ BoVW vale para todo o $d$. O que
muda é a escala. Para o mesmo $K$, VLAD e FV são 320--644$\times$ maiores do
que o BoVW. A condição D não existe (a VLAD mais pequena já tem 5120
dimensões), e o PCA é o único caminho para um descritor compacto. Comprimem-se
também melhor: $d=128$ remove 99.4--99.7\,\% da VLAD/FV $K=64$, contra
87.5\,\% do BoVW $K=1024$, retendo a mesma fração do mAP (94--97\,\%). As
dimensões são muito correlacionadas (resíduos de SIFT-PCA80 repetidos em 4
tiles). Além disso, com $\sim$900--1600 imagens de treino, o vetor de pesos do
SVM vive no espaço gerado por essas imagens (teorema do representante). Das
41\,216 dimensões do FV, no máximo $\sim$1600 direções são utilizáveis: a
dimensão nominal exagera a capacidade e não é comparável entre encodings
\cite{perronnin2010,jegou2010}.
